\documentclass[11pt]{article}
\usepackage[margin=1in]{geometry}

% Core packages
\usepackage{amsmath,amssymb}
\usepackage{array}
\usepackage{booktabs}
\usepackage{enumitem}
\usepackage{hyperref}
\usepackage{xcolor}

% Paragraphs
\setlength{\parindent}{0pt}
\setlength{\parskip}{0.8\baselineskip}
\emergencystretch=2em

% Lists
\setlist[itemize]{topsep=0.25\baselineskip,itemsep=0.25\baselineskip}
\setlist[enumerate]{topsep=0.25\baselineskip,itemsep=0.25\baselineskip}

% Table helpers
\newcolumntype{L}[1]{>{\raggedright\arraybackslash}p{#1}}

\title{Architecture Decision Record: ADR-002\\
\large Implementing Pydantic Logfire Structured Logging}
\author{Abel-Raul Mazilu}
\date{13-05-2026}

\begin{document}
\maketitle

\begin{center}
\begin{tabular}{rl}
\textbf{Status:} & Accepted \\
\textbf{Project:} & DVerse Communication Platform 2.0 \\
\end{tabular}
\end{center}

\section{Context}

The DVerse platform requires structured visibility into the behavior of its AI agents. As agents communicate over Zenoh and through the chat application, ordinary freeform logs are insufficient for debugging and analysis. Engineers need logs that can be searched, filtered, correlated with traces, and validated against consistent schemas.

The agent layer is implemented in Python and uses Pydantic models. This makes Pydantic Logfire a strong fit for structured, schema-aware logging because it can work naturally with the models already used by the system.

\section{Decision}

Sprint 4 will implement Pydantic Logfire for structured logging across the agent layer. Logs will be emitted at major system boundaries and will be designed to correlate with OpenTelemetry traces.

The implementation will add structured logs at the following boundaries:

\begin{itemize}
    \item Agent initialization and configuration loading.
    \item Incoming and outgoing Zenoh messages.
    \item Chat application events involving agents.
    \item Validation failures for agent messages or payloads.
    \item Runtime exceptions and recoverable error states.
\end{itemize}

Logs should include stable fields such as agent identifier, topic name, message identifier, trace identifier, timestamp, event type, and validation status. This will make logs queryable, consistent, and easier to correlate with OpenTelemetry traces.

\section{Alternative Decisions}

One alternative was to rely on plain JSON logs. This was rejected because plain JSON logging would require more manual serialization and would not provide the same Pydantic-native ergonomics.

Another alternative was to defer structured logging until after the communication layer was complete. This was rejected because logs are most useful when they are designed alongside the events and schemas they describe.

\section{Consequences}

\subsection{Positive}

\begin{itemize}
    \item Logs will be structured, queryable, and consistent across major agent boundaries.
    \item Pydantic model awareness will reduce manual serialization work.
    \item Logs can be correlated with OpenTelemetry traces through shared trace identifiers.
\end{itemize}

\subsection{Negative}

\begin{itemize}
    \item The team must define and maintain stable logging schemas.
    \item Logging introduces additional operational setup and runtime overhead.
    \item Sensitive fields must be reviewed carefully before being emitted to the logging system.
\end{itemize}

\section{Scope}

\begin{table}[ht]
\centering
\renewcommand{\arraystretch}{1.25}
\begin{tabular}{L{0.32\textwidth} L{0.58\textwidth}}
\toprule
\textbf{Deliverable} & \textbf{Description} \\
\midrule
Agent logs & Emit structured logs for initialization, configuration, and lifecycle events. \\
Message logs & Log incoming and outgoing Zenoh messages with stable identifiers. \\
Chat logs & Log chat application events involving AI agents. \\
Error logs & Capture validation failures, exceptions, and recoverable error states. \\
\bottomrule
\end{tabular}
\end{table}

\newpage

\section*{Statement on AI Use}

Artificial intelligence was used to support the drafting of parts of this document. The generated text was subsequently reviewed, checked, and edited where necessary by me, the author, Abel-Raul Mazilu.

\end{document}